% =====================================================================
\chapter{Soluciones de los ejercicios}
\label{ap:soluciones}
% =====================================================================

Cada solución se presenta en cinco partes: datos, planteamiento, desarrollo, resultado y comprobación. En los pasos intermedios se redondea a cuatro cifras significativas salvo que se indique otra cosa; los resultados exactos se dan como fracción cuando existe. Todos los valores numéricos se han comprobado con Python.

% ---------------------------------------------------------------------
\section{Tema 1}

\begin{solucion}{eje:1-lineal}
\datos un MLP de capas afines $z\mapsto zW_\ell+b_\ell$ sin activación.
\planteamiento aplicar la \cref{prop:colapso} y buscar un producto $W_1W_2$ de rango menor que~2.
\pasos
\begin{enumerate}[label=\arabic*.]
  \item Por la \cref{prop:colapso}, diez capas afines equivalen a una sola: $x\mapsto x(W_1\cdots W_{10})+b'$. La profundidad efectiva es~1.
  \item En una red $2\to 1\to 2$, $W_1\in\R^{2\times 1}$ y $W_2\in\R^{1\times 2}$, luego $W=W_1W_2$ es $2\times 2$ con rango a lo sumo~1.
  \item Ejemplo: $W_1=\begin{psmallmatrix}1\\0\end{psmallmatrix}$, $W_2=(1\ \ 1)$, $b_1=0$, $b_2=0$. La red calcula $x\mapsto(x_1,x_1)$. Ninguna red de esta forma calcula la identidad $x\mapsto x$, que tiene rango~2.
  \item Añadir una capa lineal sustituye $W$ por $WW_{\ell}$: el conjunto de funciones sigue siendo de la forma $x\mapsto xW+b$ (no aumenta) y el rango no puede crecer (puede disminuir).
\end{enumerate}
\resultado \textbf{profundidad efectiva~1; una capa intermedia de anchura~1 limita el rango a~1 y reduce la capacidad.}
\comprobacion con matrices $1\times 1$, $x\cdot 2\cdot 3=6x$, una sola capa; en el ejemplo, la entrada $(0,1)$ produce $(0,0)$, de modo que la información de $x_2$ se pierde.
\end{solucion}

\begin{solucion}{eje:1-xor}
\datos XOR: $(0,0)\mapsto 0$, $(0,1)\mapsto 1$, $(1,0)\mapsto 1$, $(1,1)\mapsto 0$; la red~\eqref{eq:xor}.
\planteamiento reproducir la demostración de la \cref{prop:xor} y evaluar la capa oculta en las cuatro entradas.
\pasos
\begin{enumerate}[label=\arabic*.]
  \item Un perceptrón necesitaría $b<0$, $w_1+w_2+b<0$ (ceros) y $w_2+b>0$, $w_1+b>0$ (unos).
  \item Sumando las dos primeras, $w_1+w_2+2b<0$; sumando las dos últimas, $w_1+w_2+2b>0$. Contradicción.
  \item Capa oculta: $(0,0)\mapsto(\ReLU 0,\ReLU(-1))=(0,0)$; $(0,1)\mapsto(1,0)$; $(1,0)\mapsto(1,0)$; $(1,1)\mapsto(2,1)$.
  \item En el plano $(h_1,h_2)$ los dos unos coinciden en $(1,0)$, y la recta $h_1-2h_2=\tfrac12$ separa ese punto de $(0,0)$ y $(2,1)$.
\end{enumerate}
\resultado \textbf{XOR no es linealmente separable; la capa oculta lleva las entradas a $(0,0)$, $(1,0)$, $(1,0)$, $(2,1)$, que sí lo son.}
\comprobacion $\hat y=h_1-2h_2$ da $0$, $1$, $1$ y $2-2=0$, la tabla de la \cref{fig:xor}.
\end{solucion}

\begin{solucion}{eje:1-param}
\datos MLP $784\to 256\to 128\to 10$.
\planteamiento \cref{prop:parametros-mlp}: la capa $n\to m$ tiene $(n+1)m$ parámetros.
\pasos
\begin{enumerate}[label=\arabic*.]
  \item Capa 1: $784\cdot 256+256=200\,704+256=200\,960$.
  \item Capa 2: $256\cdot 128+128=32\,768+128=32\,896$.
  \item Capa 3: $128\cdot 10+10=1\,290$.
  \item Total: $200\,960+32\,896+1\,290=235\,146$.
  \item Fracción de la primera capa: $200\,960/235\,146=0{,}8546$.
\end{enumerate}
\resultado \textbf{$235\,146$ parámetros, el $85{,}5\,\%$ en la primera capa.}
\comprobacion $785\cdot 256=200\,960$; estimación gruesa: $800\cdot 250=200\,000$.
\end{solucion}

\begin{solucion}{eje:1-cs231n}
\datos redes $3\to 4\to 2$ y $3\to 4\to 4\to 1$.
\planteamiento contar pesos $n_{\ell-1}n_\ell$ y sesgos $n_\ell$ por capa.
\pasos
\begin{enumerate}[label=\arabic*.]
  \item Primera red: pesos $3\cdot 4+4\cdot 2=12+8=20$; sesgos $4+2=6$; total $26$.
  \item Segunda red: pesos $3\cdot 4+4\cdot 4+4\cdot 1=12+16+4=32$; sesgos $4+4+1=9$; total $41$.
\end{enumerate}
\resultado \textbf{26 parámetros (20 pesos y 6 sesgos) y 41 parámetros (32 pesos y 9 sesgos).}
\comprobacion con la \cref{prop:parametros-mlp}: $(3+1)4+(4+1)2=16+10=26$ y $(3+1)4+(4+1)4+(4+1)1=16+20+5=41$.
\end{solucion}

\begin{solucion}{eje:1-or}
\datos OR: $(0,0)\mapsto 0$ y las demás $\mapsto 1$; $\eta=1$, $w=(0,0)$, $b=0$.
\planteamiento aplicar la \cref{def:regla} entrada a entrada; recuérdese que $\hat y=1$ solo si la preactivación es estrictamente positiva.
\pasos
\begin{enumerate}[label=\arabic*.]
  \item Pasada 1. $(0,0)$: $0\Rightarrow\hat y=0$, correcto. $(0,1)$: $0\Rightarrow\hat y=0$, error; $w=(0,1)$, $b=1$. $(1,0)$: $0+1=1>0$, correcto. $(1,1)$: $1+1=2$, correcto.
  \item Pasada 2. $(0,0)$: $b=1>0\Rightarrow\hat y=1$, error; $w=(0,1)$, $b=0$. $(0,1)$: $1$, correcto. $(1,0)$: $0\Rightarrow\hat y=0$, error; $w=(1,1)$, $b=1$. $(1,1)$: $3$, correcto.
  \item Pasada 3. $(0,0)$: $1>0$, error; $w=(1,1)$, $b=0$. $(0,1)$: $1$; $(1,0)$: $1$; $(1,1)$: $2$; correctos.
  \item Pasada 4. Preactivaciones $0$, $1$, $1$, $2$: sin errores.
\end{enumerate}
\resultado \textbf{4 correcciones; $w=(1,1)$, $b=0$.}
\comprobacion $\ind[x_1+x_2>0]$ vale $0$ en $(0,0)$ y $1$ en las otras tres entradas: es OR.
\end{solucion}

\begin{solucion}{eje:1-novikoff}
\datos $\theta_0=0$; actualización $\theta\leftarrow\theta+\eta s_i\tilde x_i$; AND con $\tilde x_i=(x_i,1)$ y $s=(-1,-1,-1,+1)$ para $(0,0),(0,1),(1,0),(1,1)$.
\planteamiento (a) inducción sobre el número de pasos; (b) aplicar las definiciones del \cref{teo:novikoff}.
\pasos
\begin{enumerate}[label=\arabic*.]
  \item (a) Sea $\bar\theta_k$ la sucesión obtenida con $\eta=1$. Se afirma que $\theta_k=\eta\bar\theta_k$ y que ambas ejecuciones hacen las mismas predicciones. Para $k=0$, $0=\eta\cdot 0$. Si $\theta_k=\eta\bar\theta_k$, entonces $\theta_k\cdot\tilde x=\eta\,\bar\theta_k\cdot\tilde x$ tiene el mismo signo que $\bar\theta_k\cdot\tilde x$ (y es nula a la vez), luego la predicción, el error y la actualización coinciden, y $\theta_{k+1}=\eta\bar\theta_k+\eta s_i\tilde x_i=\eta\bar\theta_{k+1}$.
  \item (b) $\norm{(1,1,-1{,}5)}=\sqrt{1+1+2{,}25}=\sqrt{4{,}25}=2{,}0616$; $u=(1,1,-1{,}5)/2{,}0616$.
  \item Márgenes sin normalizar $s_i(1,1,-1{,}5)\cdot\tilde x_i$: $(-1)(-1{,}5)=1{,}5$; $(-1)(-0{,}5)=0{,}5$; $0{,}5$; $(+1)(0{,}5)=0{,}5$. El mínimo es $0{,}5$, luego $\gamma=0{,}5/2{,}0616=0{,}2425$.
  \item $R=\max\norm{\tilde x_i}=\norm{(1,1,1)}=\sqrt3=1{,}732$.
  \item Cota: $R^{2}/\gamma^{2}=3\cdot 4{,}25/0{,}25=51$.
\end{enumerate}
\resultado \textbf{(a) las predicciones no dependen de $\eta$; (b) $\gamma\approx 0{,}2425$, $R=\sqrt3$ y $M\le 51$, compatible con las 10 correcciones observadas.}
\comprobacion la cota es una garantía en el peor caso sobre todos los órdenes de presentación; además, otro $u$ con mayor margen daría una cota menor.
\end{solucion}

% ---------------------------------------------------------------------
\section{Tema 2}

\begin{solucion}{eje:2-sigmoide}
\datos $\sigma(t)=(1+e^{-t})^{-1}$.
\planteamiento regla de la cadena sobre la potencia $-1$.
\pasos
\begin{enumerate}[label=\arabic*.]
  \item $\sigma'(t)=-(1+e^{-t})^{-2}\cdot(-e^{-t})=e^{-t}/(1+e^{-t})^{2}$.
  \item $e^{-t}/(1+e^{-t})^{2}=\frac{1}{1+e^{-t}}\cdot\frac{e^{-t}}{1+e^{-t}}$, y el segundo factor es $1-\sigma(t)$.
  \item $\sigma'(0)=\tfrac12\cdot\tfrac12=\tfrac14$.
  \item $\sigma(\ln 3)=1/(1+1/3)=3/4$, luego $\sigma'(\ln 3)=\tfrac34\cdot\tfrac14=\tfrac{3}{16}=0{,}1875$.
\end{enumerate}
\resultado \textbf{$\sigma'=\sigma(1-\sigma)$; $\sigma'(0)=1/4$; $\sigma'(\ln 3)=3/16$.}
\comprobacion la diferencia centrada con $\varepsilon=10^{-6}$ en $\ln 3$ da $0{,}1875$; el máximo de $\sigma'$ es $1/4$ (\cref{prop:sigmoide}).
\end{solucion}

\begin{solucion}{eje:2-anchura}
\datos $\FFN_{\mathrm{SwiGLU}}(x)=(\SiLU(xW_g)\odot xW_u)W_d$; MLP clásico $d\to 4d\to d$.
\planteamiento describir cada matriz e igualar el número de pesos (\cref{prop:anchura}).
\pasos
\begin{enumerate}[label=\arabic*.]
  \item $W_g\in\R^{d\times H}$: calcula la puerta, que tras SiLU decide cuánto pasa de cada componente.
  \item $W_u\in\R^{d\times H}$: calcula el contenido, que la puerta multiplica componente a componente.
  \item $W_d\in\R^{H\times d}$: proyecta el resultado de vuelta a dimensión $d$.
  \item Pesos: $3dH$ frente a $8d^{2}$; igualando, $H=8d/3=\tfrac23\cdot 4d$.
  \item Con $d=4096$: $H=10\,922{,}7$, redondeado a $11\,008=43\cdot 256$ en el código publicado de LLaMA.
\end{enumerate}
\resultado \textbf{puerta, contenido y proyección de salida; $H\approx\tfrac23\cdot 4d$ para no aumentar el número de parámetros respecto al MLP clásico.}
\comprobacion $3\cdot 4096\cdot 11\,008=135\,266\,304$ y $2\cdot 4096\cdot 16\,384=134\,217\,728$: difieren en un $0{,}8\,\%$.
\end{solucion}

\begin{solucion}{eje:2-tanh}
\datos $\tanh t=(e^{t}-e^{-t})/(e^{t}+e^{-t})$, $\sigma(t)=1/(1+e^{-t})$.
\planteamiento dividir numerador y denominador de tanh por $e^{t}$ y absorber las constantes en las capas vecinas.
\pasos
\begin{enumerate}[label=\arabic*.]
  \item $\tanh t=(1-e^{-2t})/(1+e^{-2t})$.
  \item $\tanh t+1=\bigl[(1-e^{-2t})+(1+e^{-2t})\bigr]/(1+e^{-2t})=2/(1+e^{-2t})=2\sigma(2t)$.
  \item En $t=0{,}5$: $\tanh 0{,}5+1=0{,}4621+1=1{,}4621$ y $2\sigma(1)=2\cdot 0{,}7311=1{,}4621$.
  \item Una capa tanh calcula $h=\tanh(xW+b)=2\sigma(2xW+2b)-\mathbf 1$. La capa siguiente recibe $hW'+b'=\sigma(x(2W)+2b)\,(2W')+(b'-\mathbf 1W')$.
  \item Luego la red con sigmoides de pesos $2W$, sesgos $2b$ y capa siguiente con pesos $2W'$ y sesgos $b'-\mathbf 1W'$ calcula la misma función. El recíproco se obtiene invirtiendo las sustituciones.
\end{enumerate}
\resultado \textbf{$\tanh t+1=2\sigma(2t)$; ambas activaciones generan la misma familia de funciones.}
\comprobacion en $t=0$: $\tanh 0+1=1=2\sigma(0)$.
\end{solucion}

\begin{solucion}{eje:2-swish}
\datos $s(t)=t\,\sigma(\beta t)$.
\planteamiento regla del producto y $\frac{d}{dt}\sigma(\beta t)=\beta\sigma(\beta t)(1-\sigma(\beta t))$.
\pasos
\begin{enumerate}[label=\arabic*.]
  \item $s'(t)=\sigma(\beta t)+\beta t\,\sigma(\beta t)(1-\sigma(\beta t))=\sigma(\beta t)\bigl[1+\beta t(1-\sigma(\beta t))\bigr]$.
  \item $\beta=1$, $t=1$: $0{,}7311\cdot(1+0{,}2689)=0{,}7311\cdot 1{,}2689=0{,}9277$.
  \item $\beta=1$, $t=0$: $0{,}5\cdot 1=0{,}5$.
  \item $\beta=1$, $t=-1$: $0{,}2689\cdot(1-0{,}7311)=0{,}2689^{2}=0{,}07233$.
  \item $\beta=2$, $t=1$: $\sigma(2)=0{,}8808$; $0{,}8808\cdot(1+2\cdot 0{,}1192)=0{,}8808\cdot 1{,}2384=1{,}091$.
\end{enumerate}
\resultado \textbf{$s'(t)=\sigma(\beta t)[1+\beta t(1-\sigma(\beta t))]$; con $\beta=1$: $0{,}9277$, $0{,}5$ y $0{,}07233$; con $\beta=2$ en $t=1$: $1{,}091$.}
\comprobacion con $\beta=1$ coincide con la \cref{prop:silu}(i); la pendiente $1{,}091>1$ en el último caso es imposible para la ReLU.
\end{solucion}

\begin{solucion}{eje:2-desvanece}
\datos diez capas sigmoides con $z=2$.
\planteamiento por la \cref{prop:activacion}, cada capa multiplica el gradiente por $\phi'(z)$.
\pasos
\begin{enumerate}[label=\arabic*.]
  \item $\sigma(2)=0{,}8808$; $\sigma'(2)=0{,}8808\cdot 0{,}1192=0{,}1050$.
  \item Factor: $0{,}1050^{10}=1{,}628\cdot 10^{-10}$.
  \item Cota: $4^{-10}=9{,}537\cdot 10^{-7}$; el factor real es unas $5\,900$ veces menor.
  \item Con ReLU y $z=2>0$: $\ReLU'(2)=1$ y el factor es $1^{10}=1$.
\end{enumerate}
\resultado \textbf{factor $1{,}6\cdot 10^{-10}$ con sigmoides, frente a~1 con ReLU.}
\comprobacion $\ln(0{,}1050)\cdot 10=-22{,}54$ y $e^{-22{,}54}=1{,}6\cdot 10^{-10}$.
\end{solucion}

\begin{solucion}{eje:2-swiglu}
\datos $x=(1,0)$ y las matrices del \cref{ej:swiglu}.
\planteamiento repetir los pasos del ejemplo.
\pasos
\begin{enumerate}[label=\arabic*.]
  \item $xW_g=(1\cdot 1+0\cdot 0,\ 1\cdot(-1)+0\cdot 1)=(1,-1)$.
  \item $\SiLU(1)=0{,}7311$; $\SiLU(-1)=-1\cdot\sigma(-1)=-0{,}2689$. Puerta: $(0{,}7311;\,-0{,}2689)$.
  \item $xW_u=(1\cdot 1+0\cdot 1,\ 1\cdot 0+0\cdot 1)=(1,0)$.
  \item Producto: $(0{,}7311\cdot 1;\ -0{,}2689\cdot 0)=(0{,}7311;\,0)$.
  \item Salida: $0{,}7311\cdot 1+0\cdot(-1)=0{,}7311$.
\end{enumerate}
\resultado \textbf{salida $0{,}7311$; la primera componente de la puerta vale $\SiLU(1)$ y la segunda no influye porque multiplica un contenido nulo.}
\comprobacion $\SiLU(-1)=-(1-\sigma(1))=-0{,}2689$; la coincidencia con la salida del \cref{ej:swiglu} se debe a las matrices elegidas.
\end{solucion}

\begin{solucion}{eje:2-swish-relu}
\datos $\mathrm{Swish}_\beta(t)=t\sigma(\beta t)$, $\beta>0$.
\planteamiento calcular la diferencia en $t\ge 0$ y $t<0$ y cambiar de variable $s=\beta\abs{t}$.
\pasos
\begin{enumerate}[label=\arabic*.]
  \item Si $t\ge 0$: $t\sigma(\beta t)-t=-t(1-\sigma(\beta t))=-t\,\sigma(-\beta t)$, porque $1-\sigma(u)=\sigma(-u)$.
  \item Si $t<0$: $t\sigma(\beta t)-0=-\abs{t}\,\sigma(-\beta\abs{t})$.
  \item En ambos casos $\abs{\mathrm{Swish}_\beta(t)-\ReLU(t)}=\abs{t}\,\sigma(-\beta\abs{t})=\frac1\beta\,s\,\sigma(-s)$ con $s=\beta\abs t$.
  \item Cuando $t$ recorre $\R$, $s$ recorre $[0,\infty)$; luego el supremo es $c/\beta$ con $c=\sup_{s\ge0}s\sigma(-s)$. La función $s\sigma(-s)$ es continua, vale~0 en $s=0$ y tiende a~0 en el infinito, así que el supremo se alcanza.
  \item $s\sigma(-s)=s/(1+e^{s})\le s\,e^{-s}$, y $s e^{-s}$ alcanza su máximo $1/e$ en $s=1$. Luego $c\le 1/e$.
  \item Como $c/\beta\to 0$, la convergencia es uniforme.
\end{enumerate}
\resultado \textbf{$\sup_t\abs{\mathrm{Swish}_\beta-\ReLU}=c/\beta$ con $c\approx 0{,}2785\le 1/e\approx 0{,}3679$; convergencia uniforme.}
\comprobacion numéricamente, el máximo de $s\sigma(-s)$ se alcanza en $s=1{,}2785$ y vale $0{,}2785$, que es, cambiado de signo, el mínimo de SiLU de la \cref{prop:silu}(iii), como debe ser, pues $\SiLU(-s)=-s\sigma(-s)$.
\end{solucion}

% ---------------------------------------------------------------------
\section{Tema 3}

\begin{solucion}{eje:3-existencia}
\datos el \cref{teo:au} y la \cref{prop:sierra}.
\planteamiento separar lo que garantiza el teorema (existencia) de lo que se necesita (tamaño, entrenamiento, estructura).
\pasos
\begin{enumerate}[label=\arabic*.]
  \item \emph{Tamaño.} El teorema no acota $H$; para algunas funciones la anchura necesaria crece exponencialmente, mientras que la profundidad permite reutilizar piezas (\cref{prop:sierra}).
  \item \emph{Aprendizaje y generalización.} El teorema solo asegura que la red existe; no garantiza que el descenso de gradiente la encuentre ni que generalice.
  \item \emph{Estructura.} Un modelo de lenguaje debe combinar información de muchos tokens; un MLP por token no lo hace y hace falta la atención (módulo~13).
\end{enumerate}
\resultado \textbf{existencia no es eficiencia: la profundidad gana en número de parámetros, en entrenabilidad y en estructura.}
\comprobacion LLaMA-7B apila 32 bloques con MLP de anchura $11\,008$ en lugar de una capa gigantesca.
\end{solucion}

\begin{solucion}{eje:3-codos}
\datos $g(t)=\ReLU(t)+2\ReLU(t-1)$.
\planteamiento evaluar cada término; la pendiente cambia en cada codo en el peso de la unidad correspondiente.
\pasos
\begin{enumerate}[label=\arabic*.]
  \item $g(0)=0+0=0$; $g(0{,}5)=0{,}5+0=0{,}5$; $g(1)=1+0=1$.
  \item $g(1{,}5)=1{,}5+2\cdot 0{,}5=2{,}5$; $g(2)=2+2\cdot 1=4$.
  \item Pendientes: $0$ para $t<0$; $1$ en $(0,1)$; $1+2=3$ para $t>1$.
\end{enumerate}
\resultado \textbf{$g=0$; $0{,}5$; $1$; $2{,}5$; $4$; pendientes $0$, $1$ y $3$, con codos en $t=0$ y $t=1$.}
\comprobacion $(4-1)/(2-1)=3$ y $(1-0)/(1-0)=1$.
\end{solucion}

\begin{solucion}{eje:3-cuadrado}
\datos $f(t)=t^{2}$ en $[0,1]$.
\planteamiento aplicar el paso~2 de la demostración del \cref{teo:au1d} y estudiar la diferencia entre cuerda y parábola.
\pasos
\begin{enumerate}[label=\arabic*.]
  \item (a) Nodos $0;\,0{,}25;\,0{,}5;\,0{,}75;\,1$; valores $0;\,0{,}0625;\,0{,}25;\,0{,}5625;\,1$.
  \item Pendientes $m_k=(t_{k+1}^{2}-t_k^{2})/h=t_k+t_{k+1}$: $0{,}25$; $0{,}75$; $1{,}25$; $1{,}75$.
  \item Red: $g(t)=0{,}25\ReLU(t)+0{,}5\ReLU(t-0{,}25)+0{,}5\ReLU(t-0{,}5)+0{,}5\ReLU(t-0{,}75)$.
  \item (b) En $[t_k,t_k+h]$, la cuerda es $I(t)=t_k^{2}+(2t_k+h)(t-t_k)$ y $I(t)-t^{2}=(t-t_k)(t_k+h-t)$.
  \item Ese producto de dos factores no negativos de suma $h$ es máximo en el punto medio, donde vale $(h/2)^{2}=h^{2}/4$.
  \item Con $h=1/4$: error máximo $1/64=0{,}015625$.
  \item Error menor que $10^{-4}$: $h^{2}/4<10^{-4}\iff h<0{,}02\iff n>50$. Bastan $n=51$ unidades.
\end{enumerate}
\resultado \textbf{$g$ como en el paso~3, con error máximo $1/64$; con $n$ tramos el error es $1/(4n^{2})$, y $51$ unidades dan error menor que $10^{-4}$.}
\comprobacion en $t=1/8$: $g=0{,}25\cdot 0{,}125=0{,}03125$ y $t^{2}=0{,}015625$, diferencia $0{,}015625=1/64$; con $n=2$ se recupera el $1/16$ del \cref{ej:cuadrado}.
\end{solucion}

\begin{solucion}{eje:3-sierra}
\datos $g(t)=2t$ en $[0,\tfrac12]$ y $2-2t$ en $[\tfrac12,1]$.
\planteamiento aplicar $g$ dos veces y usar la \cref{prop:sierra}.
\pasos
\begin{enumerate}[label=\arabic*.]
  \item $g$ en $0,\tfrac18,\dots,1$: $0,\tfrac14,\tfrac12,\tfrac34,1,\tfrac34,\tfrac12,\tfrac14,0$.
  \item $g^{(2)}=g\circ g$: $g(0)=0$, $g(\tfrac14)=\tfrac12$, $g(\tfrac12)=1$, $g(\tfrac34)=\tfrac12$, $g(1)=0$; resulta $0,\tfrac12,1,\tfrac12,0,\tfrac12,1,\tfrac12,0$.
  \item La gráfica de $g^{(2)}$ son dos «dientes» con picos en $t=\tfrac14$ y $t=\tfrac34$: cuatro piezas de pendiente $\pm 4$.
  \item $g^{(3)}$ tiene $2^{3}=8$ piezas; una red de una capa oculta necesita al menos $8-1=7$ unidades, frente a las $2\cdot 3=6$ de la red profunda.
\end{enumerate}
\resultado \textbf{$g^{(2)}$ vale $0,\tfrac12,1,\tfrac12,0,\tfrac12,1,\tfrac12,0$; $g^{(3)}$ tiene 8 piezas y exige al menos 7 unidades en una sola capa.}
\comprobacion la pendiente de $g^{(2)}$ en $(0,\tfrac14)$ es $(\tfrac12-0)/\tfrac18=4=2^{2}$.
\end{solucion}

\begin{solucion}{eje:3-abs-max}
\datos se buscan redes $c+\sum_ja_j\ReLU(x\cdot w_j+b_j)$.
\planteamiento usar $\abs t=\ReLU(t)+\ReLU(-t)$, $t=\ReLU(t)-\ReLU(-t)$ y $\max(t_1,t_2)=t_2+\max(t_1-t_2,0)$.
\pasos
\begin{enumerate}[label=\arabic*.]
  \item Si $t\ge 0$, $\ReLU(t)+\ReLU(-t)=t+0=\abs t$; si $t<0$, $0+(-t)=\abs t$. Red: dos unidades con pesos de entrada $+1$ y $-1$, sesgos~0 y pesos de salida $1$ y $1$.
  \item $\max(t_1,t_2)=t_2+\ReLU(t_1-t_2)$, porque si $t_1\ge t_2$ se obtiene $t_1$ y si no, $t_2$.
  \item El término lineal $t_2$ se escribe $\ReLU(t_2)-\ReLU(-t_2)$.
  \item Red: tres unidades con pesos de entrada $(1,-1)$, $(0,1)$ y $(0,-1)$, sesgos~0 y pesos de salida $1$, $1$ y $-1$.
\end{enumerate}
\resultado \textbf{$\abs t=\ReLU(t)+\ReLU(-t)$; $\max(t_1,t_2)=\ReLU(t_1-t_2)+\ReLU(t_2)-\ReLU(-t_2)$.}
\comprobacion $(3,5)$: $0+5-0=5$; $(5,3)$: $2+3-0=5$; $(-1,-4)$: $3+0-4=-1$.
\end{solucion}

\begin{solucion}{eje:3-cpwl}
\datos MLP con activaciones ReLU en las capas ocultas y salida afín.
\planteamiento inducción sobre las capas: cada unidad es afín en cada pieza, y la ReLU solo puede partir cada pieza en dos por un hiperplano.
\pasos
\begin{enumerate}[label=\arabic*.]
  \item Se llama \emph{partición poliédrica} a una división de $\R^{n}$ en un número finito de poliedros convexos (intersecciones finitas de semiespacios cerrados) que lo recubren.
  \item Hipótesis de inducción: existe una partición poliédrica $\mathcal{P}_\ell$ tal que en cada poliedro todas las componentes de $h^{(\ell)}$ son funciones afines de $x$. Para $\ell=0$, $h^{(0)}=x$ y basta $\mathcal{P}_0=\{\R^{n}\}$.
  \item En cada $Q\in\mathcal{P}_\ell$, cada preactivación $z^{(\ell+1)}_j=h^{(\ell)}W_{:,j}+b_j$ es afín, $z_j=a_j\cdot x+c_j$.
  \item Se parte $Q$ por los hiperplanos $\{a_j\cdot x+c_j=0\}$, $j=1,\dots,n_{\ell+1}$: se obtienen a lo sumo $2^{n_{\ell+1}}$ poliedros, en cada uno de los cuales cada $z_j$ tiene signo constante, y $\ReLU(z_j)$ es $z_j$ o $0$, afín en ambos casos. Así se obtiene $\mathcal{P}_{\ell+1}$, finita.
  \item La salida, afín en $h^{(L-1)}$, es afín en cada poliedro de $\mathcal{P}_{L-1}$.
  \item La continuidad se sigue de que las capas afines y la ReLU son continuas y la composición de continuas es continua.
\end{enumerate}
\resultado \textbf{toda red ReLU calcula una función continua y afín a trozos sobre una partición poliédrica finita de $\R^{n}$.}
\comprobacion en $n=1$ con una capa oculta se recupera la \cref{prop:piezas}; la red de $\max(t_1,t_2)$ del \cref{eje:3-abs-max} es afín en cada uno de los semiplanos $t_1\ge t_2$ y $t_1\le t_2$.
\end{solucion}

% ---------------------------------------------------------------------
\section{Tema 4}

\begin{solucion}{eje:4-simetria}
\datos (a) MLP con una capa oculta; (b) $p=\softmax(xW+b)$ y $\nabla_{\!W}L=x^{\top}(p-e_y)$ (\cref{prop:capa-lineal,prop:delta-salida}).
\planteamiento (a) \cref{prop:simetria}; (b) comparar los gradientes de dos columnas cuando todas son iguales.
\pasos
\begin{enumerate}[label=\arabic*.]
  \item (a) Con pesos nulos, todas las columnas de $W_1$ y todas las filas de $W_2$ son iguales; por la \cref{prop:simetria}, las unidades ocultas siguen siendo idénticas durante todo el entrenamiento y la capa equivale a una sola unidad. Con ReLU y $W_1=0$, además, $h=\ReLU(b_1)$ no depende de la entrada.
  \item (b) Si todas las columnas de $W$ son iguales y $b$ es constante, los logits son iguales y $p=(1/K,\dots,1/K)$.
  \item La columna $k$ de $\nabla_{\!W}L$ es $x^{\top}(p_k-\ind[k=y])$: vale $x^{\top}(1/K-1)$ para $k=y$ y $x^{\top}/K$ para $k\neq y$.
  \item Las columnas reciben gradientes distintos y se diferencian tras el primer paso. La etiqueta rompe la simetría porque cada salida tiene su propio objetivo.
\end{enumerate}
\resultado \textbf{(a) los pesos nulos dejan las unidades ocultas indistinguibles para siempre; (b) en la regresión softmax sí pueden inicializarse iguales, incluso a cero.}
\comprobacion la función de pérdida de la regresión softmax es convexa (módulo~8), de modo que el punto de partida no afecta al mínimo alcanzado; en un MLP, en cambio, la configuración simétrica se conserva paso a paso (\cref{prop:simetria}).
\end{solucion}

\begin{solucion}{eje:4-xavier}
\datos $\nin=512$, $\nout=2048$.
\planteamiento \cref{def:xavier} y \cref{teo:he}.
\pasos
\begin{enumerate}[label=\arabic*.]
  \item $\nin+\nout=2560$.
  \item Xavier normal: $s^{2}=2/2560=7{,}813\cdot 10^{-4}$; $s=0{,}02795$.
  \item Xavier uniforme: $a^{2}=6/2560=2{,}344\cdot 10^{-3}$; $a=0{,}04841$.
  \item He: $s^{2}=2/512=3{,}906\cdot 10^{-3}$; $s=0{,}0625$.
\end{enumerate}
\resultado \textbf{Xavier: $s\approx 0{,}0280$ y $a\approx 0{,}0484$; He: $s=0{,}0625$.}
\comprobacion por el \cref{lem:uniforme}, $a^{2}/3=2{,}344\cdot 10^{-3}/3=7{,}813\cdot 10^{-4}=s^{2}$.
\end{solucion}

\begin{solucion}{eje:4-uniforme}
\datos He: $s^{2}=2/\nin$; $\Var(U(-a,a))=a^{2}/3$.
\planteamiento igualar $a^{2}/3=2/\nin$.
\pasos
\begin{enumerate}[label=\arabic*.]
  \item $a=\sqrt{6/\nin}$; con $\nin=512$, $a=\sqrt{6/512}=\sqrt{0{,}01172}=0{,}1083$.
  \item Con $\nin=200$: $s=\sqrt{2/200}=\sqrt{0{,}01}=0{,}1$.
\end{enumerate}
\resultado \textbf{$a\approx 0{,}1083$; $s=0{,}1$.}
\comprobacion $0{,}1083^{2}/3=0{,}003906=2/512$.
\end{solucion}

\begin{solucion}{eje:4-decaimiento}
\datos ReLU, $s^{2}=1/n$, $\nin=n$; desviación típica $0{,}584$ en la capa~1.
\planteamiento \cref{teo:he}: $\Var(z^{(\ell+1)})=\tfrac12\nin s^{2}\Var(z^{(\ell)})$.
\pasos
\begin{enumerate}[label=\arabic*.]
  \item (a) $\tfrac12\cdot n\cdot\tfrac1n=\tfrac12$: la varianza se divide por~2 en cada capa.
  \item La desviación típica de $\ReLU(z)$ es proporcional a la de $z$, luego se divide por $\sqrt2$ en cada capa.
  \item (b) De la capa~1 a la~10 hay 9 pasos: $0{,}584\cdot 2^{-9/2}=0{,}584\cdot 0{,}04419=0{,}0258$.
  \item El \cref{tab:experimento} da $0{,}0254$.
  \item (c) Hace falta $\tfrac12 ns^{2}=1$, es decir, $s^{2}=2/n$ (He).
\end{enumerate}
\resultado \textbf{(a) factor $\tfrac12$; (b) predicción $0{,}0258$ frente a $0{,}0254$ medido; (c) $s^{2}=2/n$.}
\comprobacion la predicción teórica de la capa~1 es $\sqrt{\tfrac12-\tfrac1{2\pi}}=0{,}584$ (\cref{eje:4-he} con $\sigma^{2}=1$).
\end{solucion}

\begin{solucion}{eje:4-he}
\datos $z\sim\mathcal{N}(0,\sigma^{2})$ con densidad $f(t)=\frac{1}{\sigma\sqrt{2\pi}}e^{-t^{2}/(2\sigma^{2})}$.
\planteamiento \cref{lem:relu2} para el segundo momento; integral directa para el primero.
\pasos
\begin{enumerate}[label=\arabic*.]
  \item $f$ es par, luego $\E[\ReLU(z)^{2}]=\int_0^{\infty}t^{2}f(t)\,dt=\tfrac12\E[z^{2}]=\sigma^{2}/2$.
  \item $\E[\ReLU(z)]=\int_0^{\infty}t f(t)\,dt=\frac{1}{\sigma\sqrt{2\pi}}\bigl[-\sigma^{2}e^{-t^{2}/(2\sigma^{2})}\bigr]_0^{\infty}=\frac{\sigma^{2}}{\sigma\sqrt{2\pi}}=\frac{\sigma}{\sqrt{2\pi}}$.
  \item $\Var(\ReLU(z))=\frac{\sigma^{2}}{2}-\frac{\sigma^{2}}{2\pi}=\sigma^{2}\,\frac{\pi-1}{2\pi}$.
  \item Inicialización: por la \cref{prop:var-adelante}, $\Var(z^{(\ell+1)})=\nin s^{2}\cdot\sigma^{2}/2$; igualar a $\sigma^{2}$ da $s^{2}=2/\nin$.
  \item Capa~1 del experimento: entrada normal estándar, $\nin=256$, $s^{2}=2/256$; $\Var(z)=256\cdot\tfrac{2}{256}\cdot 1=2$.
  \item Desviación típica de $\ReLU(z)$: $\sqrt{2\cdot(\pi-1)/(2\pi)}=\sqrt{1-1/\pi}=\sqrt{0{,}6817}=0{,}8256$.
\end{enumerate}
\resultado \textbf{$\E[\ReLU(z)^{2}]=\sigma^{2}/2$, $\E[\ReLU(z)]=\sigma/\sqrt{2\pi}$, $\Var(\ReLU(z))=\sigma^{2}(\pi-1)/(2\pi)$; $s^{2}=2/\nin$; desviación típica $0{,}826$ en la capa~1.}
\comprobacion con $2\cdot 10^{6}$ muestras de $\mathcal{N}(0,4)$, la media de $\ReLU(z)^{2}$ es $2{,}00$, frente al valor exacto $4/2=2$; el valor $0{,}826$ coincide con el del \cref{tab:experimento}.
\end{solucion}

\begin{solucion}{eje:4-residual}
\datos $r_N=r_0+c\sum_{k=1}^{N}f_k$ con variables incorreladas de media~0.
\planteamiento la varianza de una suma de variables incorreladas es la suma de las varianzas, y $\Var(cX)=c^{2}\Var(X)$.
\pasos
\begin{enumerate}[label=\arabic*.]
  \item (a) $\Var(r_N)=\Var(r_0)+\sum_{k=1}^{N}c^{2}\Var(f_k)=v_0+Nc^{2}v$.
  \item (b) Con $c=1$: $v_0+Nv\to\infty$ cuando $N\to\infty$, si $v>0$.
  \item Con $c=1/\sqrt N$: $v_0+N\cdot\tfrac1N\cdot v=v_0+v$, independiente de $N$.
  \item (c) En un Transformer, el flujo residual acumula la salida de cada subcapa. GPT-2 multiplica por $1/\sqrt{N}$ los pesos de las capas que escriben en ese flujo~\cite[§2.3]{radford2019}, lo que divide la varianza de cada contribución por $N$ y mantiene acotada la del flujo al aumentar la profundidad.
\end{enumerate}
\resultado \textbf{$\Var(r_N)=v_0+Nc^{2}v$; con $c=1/\sqrt N$ vale $v_0+v$ para todo $N$.}
\comprobacion con $N=48$ capas residuales y $v_0=v=1$: sin escalar, $\Var=49$; escalando, $\Var=2$.
\end{solucion}

% ---------------------------------------------------------------------
\section{Tema 5}

\begin{solucion}{eje:5-softmax}
\datos $z=(\ln 4,0,0)$; clase correcta, la tercera.
\planteamiento $e^{\ln 4}=4$; $L=-\ln p_3$.
\pasos
\begin{enumerate}[label=\arabic*.]
  \item $e^{z}=(4,1,1)$, con suma~6.
  \item $p=(4/6,1/6,1/6)=(2/3,1/6,1/6)$.
  \item $L=-\ln(1/6)=\ln 6=1{,}792$ nats.
  \item Modelo uniforme: $L=\ln 3=1{,}099$ nats.
\end{enumerate}
\resultado \textbf{$p=(2/3,1/6,1/6)$; $L=\ln 6\approx 1{,}792$ nats, mayor que la pérdida uniforme $\ln 3\approx 1{,}099$.}
\comprobacion $p$ suma~1; la pérdida supera la uniforme porque el modelo concentra la probabilidad en una clase equivocada.
\end{solucion}

\begin{solucion}{eje:5-delta}
\datos $p=(0{,}5;\,0{,}3;\,0{,}2)$, $y=1$.
\planteamiento \cref{prop:delta-salida}.
\pasos
\begin{enumerate}[label=\arabic*.]
  \item $e_y=(0,1,0)$.
  \item $p-e_y=(0{,}5;\,0{,}3-1;\,0{,}2)=(0{,}5;\,-0{,}7;\,0{,}2)$.
  \item El paso de descenso resta $\eta(p-e_y)$: solo aumenta la componente con gradiente negativo, $z_1$.
\end{enumerate}
\resultado \textbf{$\nabla_{\!z}L=(0{,}5;\,-0{,}7;\,0{,}2)$; aumenta el logit de la clase correcta, $z_1$.}
\comprobacion $0{,}5-0{,}7+0{,}2=0$ (\cref{cor:suma-cero}).
\end{solucion}

\begin{solucion}{eje:5-y0}
\datos la red del \cref{ej:backprop}, ahora con $y=0$.
\planteamiento la pasada hacia delante no depende de $y$: $z^{(1)}=(2;\,0{,}5;\,-1{,}5)$, $h=(2;\,0{,}5;\,0)$, $p=(0{,}7214;\,0{,}2654;\,0{,}0132)$. Solo cambia $e_y$.
\pasos
\begin{enumerate}[label=\arabic*.]
  \item $L=-\ln 0{,}7214=0{,}3266$.
  \item $\delta^{(2)}=p-e_0=(-0{,}2786;\,0{,}2654;\,0{,}0132)$.
  \item $\nabla_{\!h}L=\delta^{(2)}W_2^{\top}$. Fila $(1,0,-1)$: $-0{,}2786-0{,}0132=-0{,}2918$. Fila $(0,2,0)$: $2\cdot 0{,}2654=0{,}5308$. Fila $(1,1,0)$: $-0{,}2786+0{,}2654=-0{,}0132$.
  \item $\delta^{(1)}=(-0{,}2918;\,0{,}5308;\,-0{,}0132)\odot(1,1,0)=(-0{,}2918;\,0{,}5308;\,0)$.
  \item $\nabla_{\!W_1}L=x^{\top}\delta^{(1)}$: fila~1 $=1\cdot\delta^{(1)}$, fila~2 $=2\cdot\delta^{(1)}$.
\end{enumerate}
\[
\nabla_{\!W_1}L=\begin{pmatrix}-0{,}2918&0{,}5308&0\\-0{,}5836&1{,}0616&0\end{pmatrix}.
\]
\resultado \textbf{$L\approx 0{,}3266$; $\delta^{(2)}=(-0{,}2786;\,0{,}2654;\,0{,}0132)$; $\nabla_{\!h}L=(-0{,}2918;\,0{,}5308;\,-0{,}0132)$; $\nabla_{\!W_1}L$ como arriba.}
\comprobacion $\delta^{(2)}$ suma~0; la pérdida es exactamente 1~nat menor que con $y=1$, porque $L_{y=1}-L_{y=0}=z_0-z_1=1$; las diferencias centradas reproducen $\nabla_{\!W_1}L$ con error menor que $10^{-9}$.
\end{solucion}

\begin{solucion}{eje:5-cs224n}
\datos $x=(1,2)$, $W_1=\begin{psmallmatrix}1&-1\\0&1\end{psmallmatrix}$, $b_1=(-1,-1)$, $W_2=\begin{psmallmatrix}2&0\\0&-2\end{psmallmatrix}$, $b_2=0$, $y=0$.
\planteamiento \cref{teo:backprop} con $\phi=\sigma$, $\phi'=\sigma(1-\sigma)=h\odot(1-h)$.
\pasos
\begin{enumerate}[label=\arabic*.]
  \item (a) $\delta^{(2)}=\hat y-e_y$; $\nabla_{\!h}J=\delta^{(2)}W_2^{\top}$; $\delta^{(1)}=\nabla_{\!h}J\odot h\odot(1-h)$; $\nabla_{\!x}J=\delta^{(1)}W_1^{\top}$.
  \item (b) $xW_1=(1,\,-1+2)=(1,1)$; $z^{(1)}=(0,0)$; $h=(0{,}5;\,0{,}5)$.
  \item $z^{(2)}=hW_2=(1,-1)$; $\hat y=(e/(e+e^{-1}),\,e^{-1}/(e+e^{-1}))=(0{,}8808;\,0{,}1192)$; $J=-\ln 0{,}8808=0{,}1269$.
  \item $\delta^{(2)}=(-0{,}1192;\,0{,}1192)$.
  \item $\nabla_{\!h}J=\delta^{(2)}W_2^{\top}=(-0{,}1192\cdot 2;\ 0{,}1192\cdot(-2))=(-0{,}2384;\,-0{,}2384)$.
  \item $h(1-h)=(0{,}25;\,0{,}25)$; $\delta^{(1)}=(-0{,}0596;\,-0{,}0596)$.
  \item $\nabla_{\!x}J=\delta^{(1)}W_1^{\top}$: con la fila $(1,-1)$ de $W_1$, $-0{,}0596+0{,}0596=0$; con la fila $(0,1)$, $-0{,}0596$.
  \item (c) $(D_x+1)H+(H+1)D_y=11\cdot 5+6\cdot 3=55+18=73$.
\end{enumerate}
\resultado \textbf{$\nabla_{\!x}J=(0;\,-0{,}0596)$; 73 parámetros.}
\comprobacion las diferencias centradas dan $(0{,}0000;\,-0{,}0596)$; además $\hat y_0=\sigma(z_0-z_1)=\sigma(2)=0{,}8808$.
\end{solucion}

\begin{solucion}{eje:5-grafo}
\datos $f(x,y)=x^{3}(x-y)$ en $(2,3)$.
\planteamiento nodos $a=x^{3}$, $b=x-y$, $f=ab$; retropropagar sumando las dos rutas de $x$ (\cref{prop:cadena-grafo}).
\pasos
\begin{enumerate}[label=\arabic*.]
  \item Grafo: $x\to a$, $x\to b$, $y\to b$, $a\to f$, $b\to f$.
  \item Hacia delante: $a=8$, $b=-1$, $f=-8$.
  \item Nodo producto: $\partial f/\partial a=b=-1$, $\partial f/\partial b=a=8$.
  \item Ruta por $a$: $\partial a/\partial x=3x^{2}=12$; contribución $-1\cdot 12=-12$.
  \item Ruta por $b$: $\partial b/\partial x=1$ y $\partial b/\partial y=-1$; contribuciones $8$ a $x$ y $-8$ a $y$.
  \item $\partial f/\partial x=-12+8=-4$; $\partial f/\partial y=-8$.
\end{enumerate}
\resultado \textbf{$f(2,3)=-8$; $\partial f/\partial x=-4$; $\partial f/\partial y=-8$.}
\comprobacion derivando directamente, $\partial f/\partial x=3x^{2}(x-y)+x^{3}=-12+8=-4$ y $\partial f/\partial y=-x^{3}=-8$.
\end{solucion}

\begin{solucion}{eje:5-suma}
\datos $\delta^{(2)}=p-e_y$, $\nabla_{\!b_2}L=\delta^{(2)}$, $\nabla_{\!W_2}L=h^{\top}\delta^{(2)}$.
\planteamiento usar el \cref{cor:suma-cero} y la linealidad del gradiente respecto a los ejemplos.
\pasos
\begin{enumerate}[label=\arabic*.]
  \item $\sum_k(\nabla_{\!b_2}L)_k=\sum_k\delta^{(2)}_k=0$.
  \item La fila $i$ de $\nabla_{\!W_2}L$ es $h_i\,\delta^{(2)}$, cuya suma es $h_i\cdot 0=0$.
  \item Interpretación: $\softmax(z+c\mathbf 1)=\softmax(z)$, luego $L(z+c\mathbf 1)=L(z)$ y la derivada direccional en la dirección $\mathbf 1$ es nula: $\nabla_{\!z}L\cdot\mathbf 1=0$.
  \item Para un lote, el gradiente es la suma (o la media) de los gradientes por ejemplo, cada uno con filas de suma nula; la suma conserva la propiedad.
\end{enumerate}
\resultado \textbf{las sumas son nulas para un ejemplo y para un lote; es consecuencia de la invariancia de la softmax por traslación.}
\comprobacion en el \cref{ej:backprop}, la fila~1 de $\nabla_{\!W_2}L$ suma $1{,}4428-1{,}4692+0{,}0264=0$.
\end{solucion}

\begin{solucion}{eje:5-swiglu}
\datos $a=xW_g$, $u=xW_u$, $m=\SiLU(a)\odot u$, $o=mW_d$; gradiente $\nabla_{\!o}L$.
\planteamiento \cref{prop:capa-lineal} para las capas lineales, regla del producto componente a componente y suma de las dos ramas que salen de $x$.
\pasos
\begin{enumerate}[label=\arabic*.]
  \item $\nabla_{\!m}L=\nabla_{\!o}L\,W_d^{\top}$ y $\nabla_{\!W_d}L=m^{\top}\nabla_{\!o}L$.
  \item $m_j=\SiLU(a_j)u_j$ depende de $a$ y $u$ solo a través de la componente $j$.
  \item $\partial m_j/\partial u_j=\SiLU(a_j)$, luego $\nabla_{\!u}L=\nabla_{\!m}L\odot\SiLU(a)$.
  \item $\partial m_j/\partial a_j=u_j\SiLU'(a_j)$, luego $\nabla_{\!a}L=\nabla_{\!m}L\odot u\odot\SiLU'(a)$, con $\SiLU'(t)=\sigma(t)(1+t(1-\sigma(t)))$ (\cref{prop:silu}).
  \item $\nabla_{\!W_g}L=x^{\top}\nabla_{\!a}L$ y $\nabla_{\!W_u}L=x^{\top}\nabla_{\!u}L$.
  \item $x$ alimenta dos ramas: $\nabla_{\!x}L=\nabla_{\!a}L\,W_g^{\top}+\nabla_{\!u}L\,W_u^{\top}$.
\end{enumerate}
\resultado \textbf{$\nabla_{\!u}L=\nabla_{\!m}L\odot\SiLU(a)$, $\nabla_{\!a}L=\nabla_{\!m}L\odot u\odot\SiLU'(a)$, $\nabla_{\!x}L=\nabla_{\!a}L\,W_g^{\top}+\nabla_{\!u}L\,W_u^{\top}$.}
\comprobacion con matrices aleatorias, las diferencias centradas con $\varepsilon=10^{-6}$ reproducen $\nabla_{\!x}L$ con error menor que $10^{-6}$; en $a=0$, $\SiLU'(0)=\tfrac12$.
\end{solucion}

% ---------------------------------------------------------------------
\section{Tema 6}

\begin{solucion}{eje:6-sobreajuste}
\datos $\PPL_{\text{ent}}=3{,}76$, $\PPL_{\text{val}}=30{,}73$.
\planteamiento $\PPL=e^{\mathcal{L}}$, luego $\mathcal{L}=\ln\PPL$.
\pasos
\begin{enumerate}[label=\arabic*.]
  \item (a) La perplejidad de validación es ocho veces la de entrenamiento y supera la del modelo uniforme ($28$): el modelo memoriza los nombres de entrenamiento y no generaliza. Es sobreajuste.
  \item (b) $\ln 3{,}76=1{,}324$ nats y $\ln 30{,}73=3{,}425$ nats.
  \item (c) Restaurar el decaimiento de pesos ($\lambda=0{,}005$ da $10{,}45$ en validación), usar más datos, reducir $d$ o $H$, o detener el entrenamiento antes (parada temprana).
\end{enumerate}
\resultado \textbf{sobreajuste; $1{,}324$ y $3{,}425$ nats; regularizar o aumentar los datos.}
\comprobacion $e^{1{,}324}=3{,}76$ y $e^{3{,}425}=30{,}7$; la referencia uniforme es $\ln 28=3{,}332$ nats, menor que $3{,}425$.
\end{solucion}

\begin{solucion}{eje:6-tabla}
\datos $N=10^{5}$ tokens, $K=10^{4}$.
\planteamiento la tabla tiene $K^{4}$ entradas; cada posición del corpus inicia a lo sumo un 4-grama.
\pasos
\begin{enumerate}[label=\arabic*.]
  \item (a) $K^{4}=(10^{4})^{4}=10^{16}$ entradas.
  \item (b) Hay $N-3=99\,997$ posiciones de inicio, luego a lo sumo $99\,997$ 4-gramas distintos.
  \item (c) Fracción: $99\,997/10^{16}\approx 10^{-11}$.
\end{enumerate}
\resultado \textbf{$10^{16}$ entradas, a lo sumo $99\,997$ observadas: una fracción del orden de $10^{-11}$.}
\comprobacion aunque todos los 4-gramas del corpus fueran distintos, quedarían sin datos $10^{16}-10^{5}\approx 10^{16}$ entradas; en el corpus de nombres la situación es análoga, con $614\,656$ casillas y $1885$ ejemplos.
\end{solucion}

\begin{solucion}{eje:6-param}
\datos $K=28$, $d=10$, $c=3$, $H=64$; total $4084$ (\cref{ej:practica}).
\planteamiento total $=Kd+(cd)H+H+HK+K$.
\pasos
\begin{enumerate}[label=\arabic*.]
  \item (a) $W_1$: $cd\cdot H=30\cdot 64=1920$.
  \item (b) Con $d=20$: $560+60\cdot 64+64+1792+28=560+3840+64+1792+28=6284$.
  \item (c) Con $c=8$: $280+80\cdot 64+64+1792+28=280+5120+64+1792+28=7284$.
\end{enumerate}
\resultado \textbf{(a) $1920$; (b) $6284$; (c) $7284$.}
\comprobacion ampliar el contexto de 3 a 8 solo añade $5\cdot 10\cdot 64=3200$ parámetros, crecimiento lineal; una tabla pasaría de $28^{4}$ a $28^{9}\approx 1{,}06\cdot 10^{13}$ entradas.
\end{solucion}

\begin{solucion}{eje:6-embgrad}
\datos contexto $(\texttt{a},\texttt{n},\texttt{a})$, $d=2$, $\nabla_{\!e}L=(0{,}1;\,-0{,}2;\,0{,}3;\,0;\,-0{,}4;\,0{,}5)$.
\planteamiento \cref{prop:grad-embedding}: partir $\nabla_{\!e}L$ en bloques de $d=2$ y sumar los de cada token.
\pasos
\begin{enumerate}[label=\arabic*.]
  \item Bloques: $g_1=(0{,}1;\,-0{,}2)$ (\texttt{a}), $g_2=(0{,}3;\,0)$ (\texttt{n}), $g_3=(-0{,}4;\,0{,}5)$ (\texttt{a}).
  \item $\nabla_{\!C[\texttt{a}]}L=g_1+g_3=(0{,}1-0{,}4;\,-0{,}2+0{,}5)=(-0{,}3;\,0{,}3)$.
  \item $\nabla_{\!C[\texttt{n}]}L=g_2=(0{,}3;\,0)$.
  \item $\nabla_{\!C[\texttt{e}]}L=(0,0)$: \texttt{e} no aparece en el contexto.
\end{enumerate}
\resultado \textbf{$(-0{,}3;\,0{,}3)$, $(0{,}3;\,0)$ y $(0,0)$.}
\comprobacion la suma de todas las filas de $\nabla_{\!C}L$, $(0,0{,}3)$, coincide con la suma de los tres bloques, $(0{,}1+0{,}3-0{,}4;\,-0{,}2+0+0{,}5)=(0;\,0{,}3)$.
\end{solucion}

\begin{solucion}{eje:6-bigrama}
\datos logits $y(u)=C[u]W+b$ con $C\in\R^{K\times d}$, $W\in\R^{d\times K}$, $b\in\R^{1\times K}$; modelo de bigramas $P(v\mid u)>0$.
\planteamiento (a) construir $C$, $W$, $b$ explícitos; (b) escribir la matriz de logits como producto más una matriz de rango~1 y usar la invariancia de la softmax por traslación.
\pasos
\begin{enumerate}[label=\arabic*.]
  \item (a) Si $d\ge K$, tómese $C=(I_K\ \ 0)\in\R^{K\times d}$, $W$ con sus $K$ primeras filas iguales a la matriz $A_{uv}=\ln P(v\mid u)$ y el resto nulas, y $b=0$. Entonces $y(u)=A_{u,:}$ y $\softmax(y(u))_v=P(v\mid u)/\sum_{v'}P(v'\mid u)=P(v\mid u)$.
  \item (b) La matriz de logits es $Z=CW+\mathbf 1b$. Como $\operatorname{rg}(CW)\le d$ y $\operatorname{rg}(\mathbf 1b)\le 1$, $\operatorname{rg}Z\le d+1$ (módulo~3).
  \item Las log-probabilidades son $\ln P(v\mid u)=Z_{uv}-\LSE(Z_{u,:})$, es decir, $A=Z-s\mathbf 1^{\top}$ con $s_u=\LSE(Z_{u,:})$; luego $\operatorname{rg}A\le d+2$.
  \item Para ver que hay modelos no representables, sea $P(v\mid u)\propto e^{\kappa\ind[u=v]}$ con $\kappa>0$. Entonces $A=\kappa I-s_0J$, con $s_0=\ln(e^{\kappa}+K-1)$ y $J$ la matriz de unos.
  \item Los valores propios de $\kappa I-s_0J$ son $\kappa$ (multiplicidad $K-1$) y $\kappa-Ks_0$, que es negativo porque $s_0>\kappa$. Ninguno es nulo y $\operatorname{rg}A=K$.
  \item Si $d+2<K$, el paso~3 impide representar este modelo.
\end{enumerate}
\resultado \textbf{con $d\ge K$ se representa cualquier modelo de bigramas positivo; con $d+2<K$ el rango de las log-probabilidades está acotado por $d+2<K$ y hay modelos que no se pueden representar.}
\comprobacion con $K=3$ y $\kappa=\ln 2$: $s_0=\ln 4$ y los valores propios son $\ln 2$ y $\ln 2-3\ln 4=-5\ln 2\neq 0$. La limitación se conoce como \emph{cuello de botella de la softmax}~\cite{yang}.
\end{solucion}
